% 12 Apr 2026: Clean up

\section{Twenty-Second Sunday after Trinity: He who has ears to hear, let him hear}

\begin{quote}

Mark 4:21-25. And He said to them: Does the light come in to be put under the bushel or under the bench? Is it not rather that it should be set upon the lampstand? For nothing is hidden except that it should be made manifest, nor has anything come to pass in order to remain hidden, but that it should come to light. If anyone has ears to hear, let him hear! And He said to them: Take heed what you hear! With the measure you measure, it shall be measured to you, and to you who hear there shall be given yet more. For he who has, to him shall be given; and he who has not, from him shall be taken even that which he has.

\end{quote}

\bigskip

Hearing and speaking are so common that they hardly seem worth noticing. Nevertheless, both are of such importance that not only life and death in the purely natural sense often depend upon them, but eternal life and eternal death also depend upon the Word and on rightly hearing it. For there is a mutual connection between the two.

Just as the Word is called a two-edged sword, revealing the heart’s secret counsels, and that for some it becomes a smell of death leading to death, and for others a smell of life leading to life, so it is also said of hearing: "Faith comes from hearing, but hearing comes through the Word of God" (Rom. 10:17). Therefore the Lord says in our text: "Take heed what you hear!" And in another place: "See therefore how you hear!" (Luke 8:18).

A man’s salvation truly depends both on what he hears and on how he hears.

No one who can hear can avoid hearing what is said all around him. But it is not long before he begins to make distinctions and to prefer hearing one thing rather than another.

The ancients said: "Tell me with whom you keep company, and I shall tell you who you are." One could just as well recast it thus: "Tell me what you like to hear, and I shall tell you who you are."

Some like to be where the talk is empty and frivolous, the sort of talk a dog would be ashamed of, if it were written out and read aloud again. They do so for pastime, they say, and readily add some sort of diversion or game besides.

Others prefer pleasant company where they can hear gossip, evil report, and slander of their neighbors, and they drink in such talk as the most delightful draught.

Or they gather with those who let foul speech proceed out of their mouth together with mocking and frivolous jesting, immoral songs, oaths, and blasphemies, and their souls are permeated with venom and poison from hell.

Some move in far more elevated spheres. They care only for hearing lofty subjects discussed: science, music, art, poetry, and all such things as require so costly an education that only a few can hope to gain entrance into that esteemed society. And accordingly they find their pleasure and delight in such things.

... all these different ways of listening leave their own decisive mark on the soul; but none of them brings peace. And however much anyone may comfort himself with being a Christian, though he has his delight in lending his ear to one or more of the above-mentioned ways of listening, and thinks he can do so unharmed, he deceives himself. No one can withdraw himself from the influence of that to which he listens with delight and gladness; daily experience convinces us of it.

Only one thing truly does the soul good to hear: it requires neither art nor learning to understand, which need not hide itself “under the bushel or under the bench” either; it is that which, like a light on a lampstand and a city on a hill, shines for all the world and for every single soul with such heavenly power that it penetrates into the soul’s inmost depths and reveals all that is concealed: the most secret counsels and thoughts, the whole nature and abomination of sin; now it sounds like a trumpet of judgment, proclaiming the fire that is never quenched and the worm that never dies, now like a still, loving voice from heaven: “Even if a mother forgets her nursing child, I have not forgotten you!”—it is that by which the thorn-crowned and ascended Savior Himself is made manifest: the Word, the Word of God.

"Hear the Word of the Lord, you princes of Sodom," cries the Prophet; "hear My voice, and I will be your God, and you shall be My people." "Your words were found by me," says Jeremiah, "and I ate them, and Your words were to me joy and the gladness of my heart."

\textbf{Let all the earth hear the Word of the Lord!}

Friend, you have ears; do you hear the Word of God? Then blessed are you, if you act on it—that is, if you hear rightly.

For, says the Lord, "See therefore how you hear."

The Word does not return empty. Where the Gospel is heard, it works either life and blessedness or judgment and death; it is—as it is said of the Son of God—appointed for the fall and rising of many in Israel, and for a sign that is opposed. He who hears the Word and acts accordingly builds his house upon a rock; but he who does not act accordingly sets his house upon sand, and its fall shall be great.

A great and holy responsibility therefore rests upon every one who hears the Word of God, and to whom this precious treasure of life is thus entrusted.

To Israel the Lord says: "From the day your fathers came out of the land of Egypt to this day, I sent to you all My servants the prophets, daily, early, and continually. But they did not hear Me, and did not bend their ear to Me; they hardened their neck, they acted worse than their fathers. Therefore I will cast you out from before My face, just as I cast out your brethren, all the offspring of Ephraim."

Does not that same Word and that same judgment still sound to us?
"With the measure you measure, it shall be measured to you again, and to you who hear there shall be given yet more—but from him who has not there shall be taken even that which he has."

In these words the Lord exhorts us to deal carefully with the Word of God, and to be readier to hear than a fool is to offer sacrifice.

If you go to hear or read the Word of God with prayer, crying out for the Spirit’s blessing, and with a hungry and thirsty heart, so that you think you could gladly swallow the whole of it in order to satisfy your soul, then the Lord shall give you according to your longing and desire, "a good measure, heaped up, full, and overflowing, in your bosom." "In joy you shall go out and in peace you shall be led forth; the mountains and the hills shall break forth into shouts of joy before your face, and all the trees of the field shall clap their hands" (Isa. 55:12).

But if you go to the Word with indifference, with the proud thought that you yourself can master it, without hunger and thirst, for mere show or outward appearance, then your heart will grow emptier and emptier, weaker and weaker, until the words spring back from your soul like a projectile hurled against a wall, until you become hard as stone, and judgment comes upon you unawares as upon Israel. Then even that which you had is taken from you.

And now, friend, how often have you been in God’s house? How often have you listened to His Word of sin and grace, of death and life? How lovingly has the Lord called to you, how gently and earnestly has the Father drawn you again and again—and how do you hear? Is it with delight and longing, are the words sweeter to you than honey and honeycomb, do you draw water with joy out of the springs of life, or are you sluggish and dull, is it only reluctantly or by habit that you take in the Savior’s message, is it a relief to you when the preaching of the Word is over, do you hasten with double eagerness to your newspaper, your neighbor’s talk, or to your books and speculations?

Brothers and sisters! Do you remember that the Lord has spoken to you daily, early, and continually; that the hourglass is nearly run out; that you must answer for all the words God has spoken to you in vain; and that judgment follows hard upon hardening?

He who has ears to hear, let him hear.




